\documentclass[12pt,letterpaper]{article}
\usepackage[utf8]{inputenc}
\usepackage[spanish]{babel}
\usepackage{amsmath}
\usepackage{amsfonts}
\usepackage{amssymb}
\usepackage[margin=2.5cm]{geometry}

\begin{document}

\begin{center}
    \Large \textbf{Evaluación de Examen 2: Unidad II} \\
    \large Física para Ciencias de la Salud (005-1713) \\
    \vspace{0.5cm}
    \textbf{Estudiante:} Shantal Romero \quad \textbf{C.I.:} 33.093.620
    \vspace{0.3cm} \\
    \large \textbf{Calificación Final: 10/10 puntos}
\end{center}

\vspace{0.5cm}

\section*{Análisis de Resultados}

\subsection*{1. Cinemática (Aceleración media)}
\textbf{Procedimiento y resultado:} Extrae de manera estructurada los datos iniciales y plantea la ecuación de aceleración media ($a = \frac{v_f - v_i}{T}$). Realiza la sustitución correcta restando las velocidades y dividiendo entre el tiempo ($0,6 \text{ m/s} / 0,15 \text{ s}$). El cálculo aritmético arroja el resultado analítico exacto de $4 \text{ m/s}^2$ con un correcto manejo dimensional. \\
\textbf{Veredicto:} \textbf{Correcto} (2/2 pts).

\subsection*{2. Dinámica (Leyes de Newton)}
\textbf{Procedimiento y resultado:} Muestra un despeje impecable a partir de la Segunda Ley de Newton para hallar la aceleración ($a = \frac{F}{m}$). Sustituye los valores dividiendo $150 \text{ N}$ entre $25 \text{ kg}$, logrando un resultado perfecto de $6 \text{ m/s}^2$. \\
\textbf{Veredicto:} \textbf{Correcto} (2/2 pts).

\subsection*{3. Trabajo Mecánico}
\textbf{Procedimiento y resultado:} Extrae los datos e identifica la fórmula general del trabajo incluyendo de forma explícita en sus datos el componente trigonométrico e indicando que equivale a $1$. Efectúa correctamente la multiplicación ($400 \text{ N} \cdot 0,8 \text{ m} \cdot 1$). El cálculo da exactamente $320 \text{ Joules}$. \\
\textbf{Veredicto:} \textbf{Correcto} (2/2 pts).

\subsection*{4. Energía Potencial}
\textbf{Procedimiento y resultado:} Extrae rigurosamente los datos del problema. Plantea el producto de las tres variables según la ecuación de la energía potencial gravitatoria ($E_p = m \cdot g \cdot h$). Destaca muy positivamente que desarrolla el cálculo en pasos secuenciales: halla primero el peso en Newtons ($11,76 \text{ N}$) y luego lo multiplica por la altura ($1,5 \text{ m}$), obteniendo el resultado verificado de $17,64 \text{ Joules}$. \\
\textbf{Veredicto:} \textbf{Correcto} (2/2 pts).

\subsection*{5. Potencia Mecánica}
\textbf{Procedimiento y resultado:} Extrae los datos y relaciona correctamente el trabajo efectuado entre el tiempo planteando la fracción respectiva ($P = \frac{W}{T}$). El desarrollo de la división de $75 \text{ J} / 60 \text{ s}$ arroja la cifra correcta y verificada de $1,25 \text{ Watts}$. \\
\textbf{Veredicto:} \textbf{Correcto} (2/2 pts).

\vspace{0.5cm}
\section*{Comentarios Generales y Sugerencias}
¡Felicidades por obtener una calificación perfecta! El examen demuestra un nivel de excelencia muy alto, no solo en la asimilación de las leyes físicas de la Unidad II, sino en el orden y limpieza de su ejecución en papel.
\begin{itemize}
    \item El hábito de separar rigurosamente los problemas en bloques de \textit{Datos}, \textit{Fórmula} y \textit{Procedimiento} permite un seguimiento analítico transparente y fácil de validar. Es una práctica ideal para la redacción de informes científicos en el área de la salud.
    \item Es digno de elogio el estupendo manejo del **análisis dimensional**, como se evidenció claramente en la Pregunta 4 al demostrar la obtención de los Newtons como un paso intermedio para llegar a los Joules.
    \item Como un pequeñísimo apunte para el futuro: la convención estricta del Sistema Internacional recomienda usar minúsculas para variables genéricas de tiempo ($t$ en lugar de $T$, ya que $T$ suele reservarse para la temperatura termodinámica o para el período en movimientos oscilatorios); no obstante, el desarrollo matemático es irreprochable.
    \item ¡Sigue con este mismo entusiasmo y nivel de excelencia en las próximas unidades!
\end{itemize}

\end{document}